% Unit 4 exam -- 20 MCQ (no calculator) + 1 long FRQ. 52 min, 30 points.
\documentclass{apchem-exam}

\apcunit{4}
\apcunittitle{Chemical Reactions}
\apcdoctitle{Unit Exam}
\apcsubtitle{Mon Oct 26 \textbullet\ 52 minutes \textbullet\ periodic table and equations sheet provided}

\begin{document}
\apctitleblock

\examsection{I}{Multiple Choice}{20 questions \textbullet\ 30 minutes \textbullet\ 1 point each}{\nocalculator}

% ---- 4.4 physical vs chemical (3) ------------------------------------------
\begin{mcq}
  Which process is a chemical change?
  \choices
    \choice{ice melting}
    \choice{sugar dissolving in water}
    \correct{\ce{HCl} gas dissolving in water}
    \choice{dry ice subliming}
\end{mcq}
\why{\ce{HCl} ionizes to \ce{H+} and \ce{Cl-} --- species that did not exist
before. The others leave the substance intact.}
\distractor{B}{sugar disperses as whole molecules --- physical}

\begin{mcq}
  A student observes bubbling when a solid is added to a warm liquid. This
  observation alone
  \choices
    \choice{proves a chemical reaction occurred}
    \correct{is consistent with, but does not prove, a chemical reaction}
    \choice{proves no reaction occurred}
    \choice{indicates a physical change only}
\end{mcq}
\why{Bubbles can be dissolved air escaping on warming.}

\begin{mcq}
  Melting a molecular solid requires overcoming
  \choices
    \choice{covalent bonds within molecules}
    \correct{intermolecular forces between molecules}
    \choice{ionic bonds}
    \choice{metallic bonds}
\end{mcq}
\why{Identity is unchanged, so only IMFs break --- a physical change.}

% ---- 4.3 representations (3) -----------------------------------------------
\begin{mcq}
  A correct particulate diagram of aqueous \ce{Na2SO4} shows
  \choices
    \choice{intact \ce{Na2SO4} units}
    \correct{twice as many \ce{Na+} particles as \ce{SO4^2-} particles, all
      separated}
    \choice{equal numbers of separated \ce{Na+} and \ce{SO4^2-}}
    \choice{separated \ce{Na+}, \ce{S^2-}, and \ce{O^2-}}
\end{mcq}
\why{Complete dissociation in the 2:1 ratio; the polyatomic ion stays
intact.}
\distractor{D}{polyatomic ions do not break apart on dissolving}

\begin{mcq}
  In a particulate ``before and after'' pair of diagrams, which quantity
  must be identical in both boxes?
  \choices
    \choice{the number of molecules}
    \correct{the number of atoms of each element}
    \choice{the number of distinct substances}
    \choice{the volume occupied}
\end{mcq}
\why{Atoms are conserved; molecule counts routinely change.}

\begin{mcq}
  In a particulate diagram of a completed precipitation reaction, the
  precipitate should be drawn as
  \choices
    \choice{separated ions dispersed in solution}
    \correct{a cluster of ions held together}
    \choice{neutral molecules}
    \choice{omitted entirely}
\end{mcq}
\why{A precipitate is an insoluble lattice, not free ions.}

% ---- 4.2 net ionic (4) -----------------------------------------------------
\begin{mcq}
  Which species should be written as separated ions in a complete ionic
  equation?
  \choices
    \choice{\ce{H2O(l)}}
    \choice{\ce{HC2H3O2(aq)}}
    \correct{\ce{KNO3(aq)}}
    \choice{\ce{CaCO3(s)}}
\end{mcq}
\why{Only strong electrolytes are split; \ce{KNO3} is a soluble salt.}
\distractor{B}{a weak acid stays intact}

\begin{mcq}
  The net ionic equation for mixing \ce{Na2CO3(aq)} and \ce{CaCl2(aq)} is
  \choices
    \choice{\ce{Na+(aq) + Cl^-(aq) -> NaCl(s)}}
    \correct{\ce{Ca^2+(aq) + CO3^2-(aq) -> CaCO3(s)}}
    \choice{\ce{Na2CO3(aq) + CaCl2(aq) -> CaCO3(s) + 2 NaCl(aq)}}
    \choice{no reaction occurs}
\end{mcq}
\why{Calcium carbonate is insoluble; \ce{Na+} and \ce{Cl-} are spectators.}
\distractor{C}{that is the formula equation}

\begin{mcq}
  Which pair of reactants gives the net ionic equation
  \ce{H+(aq) + OH^-(aq) -> H2O(l)}?
  \choices
    \correct{\ce{HNO3(aq)} and \ce{KOH(aq)}}
    \choice{\ce{HF(aq)} and \ce{NaOH(aq)}}
    \choice{\ce{HCl(aq)} and \ce{NH3(aq)}}
    \choice{\ce{HC2H3O2(aq)} and \ce{KOH(aq)}}
\end{mcq}
\why{Requires a strong acid AND a strong base; the others involve a weak
species that stays intact.}

% ---- 4.7 reaction types / solubility (4) -----------------------------------
\begin{mcq}
  Mixing which pair produces a precipitate?
  \choices
    \choice{\ce{KNO3(aq)} and \ce{NaCl(aq)}}
    \correct{\ce{AgNO3(aq)} and \ce{K2CrO4(aq)}}
    \choice{\ce{NH4Cl(aq)} and \ce{NaNO3(aq)}}
    \choice{\ce{LiBr(aq)} and \ce{KNO3(aq)}}
\end{mcq}
\why{\ce{Ag2CrO4} is insoluble; every other combination is soluble.}

\begin{mcq}
  The reaction \ce{Zn(s) + CuSO4(aq) -> ZnSO4(aq) + Cu(s)} is best
  classified as
  \choices
    \choice{precipitation}
    \choice{acid--base}
    \correct{oxidation--reduction}
    \choice{decomposition}
\end{mcq}
\why{Zn goes $0 \to +2$ and Cu goes $+2 \to 0$.}

\begin{mcq}
  Two clear solutions are mixed and no precipitate forms. The best
  conclusion is that
  \choices
    \choice{no ions are present}
    \correct{all possible ion combinations are soluble}
    \choice{the solutions were too dilute to react}
    \choice{a gas must have formed instead}
\end{mcq}
\why{``No reaction'' is a legitimate outcome when every product is soluble.}

% ---- 4.8 acid-base (3) -----------------------------------------------------
\begin{mcq}
  Which is a weak acid?
  \choices
    \choice{\ce{HCl}}
    \choice{\ce{HNO3}}
    \correct{\ce{HF}}
    \choice{\ce{HClO4}}
\end{mcq}
\why{HF is weak despite fluorine's electronegativity --- the H--F bond is
strong.}

\begin{mcq}
  A Bronsted--Lowry base is defined as a substance that
  \choices
    \choice{produces hydroxide ions in water}
    \correct{accepts a proton}
    \choice{donates a proton}
    \choice{dissolves to give a conducting solution}
\end{mcq}
\why{This broader definition covers \ce{NH3}.}

\begin{mcq}
  In the net ionic equation for \ce{HF(aq)} reacting with \ce{NaOH(aq)},
  hydrofluoric acid appears as
  \choices
    \correct{\ce{HF(aq)}, an intact molecule}
    \choice{\ce{H+(aq)} and \ce{F^-(aq)}}
    \choice{\ce{H+(aq)} only}
    \choice{a spectator}
\end{mcq}
\why{Weak acids are barely ionized, so the molecular form is what reacts.}

% ---- 4.9 redox (4) ---------------------------------------------------------
\begin{mcq}
  The oxidation state of nitrogen in \ce{NO2-} is
  \choices
    \choice{$+5$}
    \correct{$+3$}
    \choice{$-3$}
    \choice{$+4$}
\end{mcq}
\why{$x + 2(-2) = -1 \Rightarrow x = +3$.}
\distractor{A}{that is nitrogen in nitrate, \ce{NO3-}}

\begin{mcq}
  In \ce{2 Mg(s) + O2(g) -> 2 MgO(s)}, magnesium is
  \choices
    \correct{oxidized and acts as the reducing agent}
    \choice{oxidized and acts as the oxidizing agent}
    \choice{reduced and acts as the reducing agent}
    \choice{reduced and acts as the oxidizing agent}
\end{mcq}
\why{Mg goes $0 \to +2$ (loses electrons); by donating them it reduces the
oxygen.}
\distractor{B}{agent labels swapped --- the oxidized species is the reducing
agent}

\begin{mcq}
  Which reaction is NOT a redox reaction?
  \choices
    \choice{\ce{2 Na + Cl2 -> 2 NaCl}}
    \choice{\ce{CH4 + 2 O2 -> CO2 + 2 H2O}}
    \correct{\ce{AgNO3(aq) + NaCl(aq) -> AgCl(s) + NaNO3(aq)}}
    \choice{\ce{Fe2O3 + 3 CO -> 2 Fe + 3 CO2}}
\end{mcq}
\why{In the precipitation reaction every element keeps its oxidation state.}

% ---- 4.5/4.6 stoichiometry and titration (2) -------------------------------
\begin{mcq}
  \SI{25.0}{\milli\liter} of \ce{HCl} is exactly neutralized by
  \SI{25.0}{\milli\liter} of \SI{0.100}{\Molar} \ce{NaOH}. The
  concentration of the acid is
  \choices
    \choice{\SI{0.0500}{\Molar}}
    \correct{\SI{0.100}{\Molar}}
    \choice{\SI{0.200}{\Molar}}
    \choice{cannot be determined}
\end{mcq}
\why{Equal volumes, 1:1 stoichiometry, so equal concentrations.}

\begin{mcq}
  Titrating \ce{H2SO4} with \ce{NaOH} and assuming a 1:1 mole ratio will
  make the calculated acid concentration
  \choices
    \correct{twice the true value}
    \choice{half the true value}
    \choice{unchanged}
    \choice{four times the true value}
\end{mcq}
\why{\ce{H2SO4} is diprotic, so the true ratio is 1 acid : 2 base;
attributing all the base to a 1:1 reaction doubles the apparent acid.}

\examsection{II}{Free Response}{1 question \textbullet\ 20 minutes}{\calculatorok}

\frq{long}{10}{CED 4.2--4.9 synthesis \textbullet\ SP 1, 4, 6}

A student mixes \SI{40.0}{\milli\liter} of \SI{0.150}{\Molar}
\ce{BaCl2(aq)} with \SI{40.0}{\milli\liter} of \SI{0.100}{\Molar}
\ce{Na2SO4(aq)}. A white solid forms.

\begin{parts}
  \item Write the balanced formula equation, including states.
        \answerlines[2]{\ce{BaCl2(aq) + Na2SO4(aq) -> BaSO4(s) +
        2 NaCl(aq)}}
  \item Write the net ionic equation and identify the spectator ions.
        \answerlines[2]{\ce{Ba^2+(aq) + SO4^2-(aq) -> BaSO4(s)};
        spectators \ce{Na+} and \ce{Cl-}.}
  \item Determine the limiting reactant. Show the comparison.
        \workspace[2.6cm]
        \keyonly{\begin{apcanswer}$n(\ce{Ba^2+}) = (0.0400)(0.150) =
        \SI{6.00e-3}{\mole}$; $n(\ce{SO4^2-}) = (0.0400)(0.100) =
        \SI{4.00e-3}{\mole}$. The ratio is 1:1, so sulfate is
        limiting.\end{apcanswer}}
  \item Calculate the mass of \ce{BaSO4}
        ($M = \SI{233.39}{\gram\per\mole}$) produced.
        \answerlines[2]{$4.00\times10^{-3} \times 233.39 =
        \SI{0.934}{\gram}$}
  \item Calculate the concentration of the excess \ce{Ba^2+} remaining after
        the reaction. Assume volumes are additive. \workspace[2.4cm]
        \keyonly{\begin{apcanswer}Excess $= 6.00\times10^{-3} -
        4.00\times10^{-3} = \SI{2.00e-3}{\mole}$ in
        \SI{80.0}{\milli\liter}: $[\ce{Ba^2+}] = 2.00\times10^{-3}/0.0800 =
        \SI{0.0250}{\Molar}$\end{apcanswer}}
  \item Explain why the formation of a white solid is stronger evidence of
        a chemical reaction than a temperature change would be.
        \answerlines[3]{A solid appearing from two clear solutions requires
        new bonding --- no physical process converts two liquids into an
        insoluble lattice. A temperature change also accompanies simple
        dissolving, so it is ambiguous on its own.}
\end{parts}

\begin{rubric}
  \rpoint{1}{(a) balanced equation with correct states}
  \rpoint{2}{(b) 1 pt net ionic equation; 1 pt both spectator ions}
  \rpoint{2}{(c) 1 pt both mole values computed; 1 pt names \ce{SO4^2-} as
    limiting with the 1:1 ratio cited}
  \rpoint{1}{(d) \SI{0.934}{\gram} with units}
  \rpoint{2}{(e) 1 pt excess moles by subtraction; 1 pt divides by the
    TOTAL \SI{80.0}{\milli\liter} --- using \SI{40.0}{\milli\liter} earns 0}
  \rpoint{2}{(f) 1 pt identifies solid formation as requiring new bonding;
    1 pt notes temperature change is ambiguous (also occurs on dissolving)}
\end{rubric}

\examtotal

\end{document}
